\chapter{Three formats, and which you are in}
\label{ch:formats}

\section*{What this chapter is for}

The round where you build something comes in three shapes, they are scored
differently, and the most common preparation error is to assume one of them.
This chapter is about finding out which you are in, and about what the three
share so that preparing for all of them does not cost three times as much.

\section{The three}

\textbf{The AI-assisted live build.} One to three hours, shared screen, you
choose your tooling, and how you use the model is explicitly part of what is
being assessed. The clearest published description of the shape runs in three
phases \dash{} \reported{a planning session in which the candidate drives the
definition of what to build, then a two-hour build with free choice of tooling
and freedom to change scope, then a review in which judgement and product
thinking are graded alongside what actually
works.}{sierra-ai-native-interview}

\textbf{The take-home.} \reported{Two to three hours, present in about a third
of disclosed loops, and usually replacing an on-site round rather than adding
one.}{field-guide-process} Nobody is watching, which changes everything:
Chapter~\ref{ch:takehome} is about the craft that replaces the watching.

\textbf{The AI-free round.} Still common, and not going away.
\reported{Several employers ban assistance explicitly, and in-person
interviewing has been rising rather than falling.}{field-guide-trends} Here the
assessment is closer to the conventional one, and the AI-specific preparation
does not apply.

\section{The split is genuine, so ask}

There is no default to assume. \reported{One industry survey of hiring leaders
reports roughly half permitting assistance, broadly or with constraints, a
fifth deciding case by case, and the remainder not permitting
it.}{coderpad-hiring-2026}

Among those that permit it, the rules differ in ways that matter.
\reported[single]{At least one large employer's round makes the assistant
available in a chat panel but does not let it edit files, and scores prompt
quality, output validation and debugging explicitly.}{google-code-comprehension}
\reported[single]{At the same employer family, reporting suggests that at more
senior levels the assisted round replaces the unassisted one rather than
joining it.}{meta-ai-enabled-coding}

\begin{kitnote}
Both of those sentences say \emph{one source} because they rest on one, and the
second is the more decision-relevant of the two. It is exactly the kind of
specific, confident, widely-repeated claim that this book's Appendix~\ref{app:refused}
refuses when it cannot be traced. It is included here, marked, because it
changes what you would prepare \dash{} but do not plan a week around it without
asking the company directly.
\end{kitnote}

\begin{kitmethod}
Ask in the screen, in one message, and ask all four:

\textbf{1.} Is there a build round, and is it live or a take-home?

\textbf{2.} Is AI assistance permitted? If so, is it a specific tool, or my own
choice?

\textbf{3.} How long, and do you expect something running at the end?

\textbf{4.} Is there a brief in advance?

These are administrative questions with administrative answers. Not asking them
is the expensive choice.
\end{kitmethod}

\section{What the three have in common}

Most of Part~\ref{part:build} applies to all three, which is why preparing for
all of them is cheaper than it looks.

\begin{kittable}{@{}lccc@{}}
& \textbf{Live AI} & \textbf{Take-home} & \textbf{AI-free} \\
\midrule
Discovery from a vague brief (Ch.~\ref{ch:discovery}) & yes & yes & yes \\
Design that can be drawn (Ch.~\ref{ch:design}) & yes & written & yes \\
Time-boxing with fallbacks (Ch.~\ref{ch:runsheet}) & yes & yes & yes \\
Vertical slices (Ch.~\ref{ch:slices}) & yes & yes & yes \\
Visible AI use (Ch.~\ref{ch:ai-visibly}) & \textbf{scored} & partly & no \\
Decomposition (Ch.~\ref{ch:decomposition}) & yes & yes & partly \\
Failure handling (Ch.~\ref{ch:breaks}) & yes & yes & yes \\
\end{kittable}

One row differs, and Chapter~\ref{ch:ai-visibly} is that row.

\kitfig{three-formats}{% Chapter 19. One question in the screen decides which column you are in.
\begin{tikzpicture}[node distance=7mm]
  \node[kitbox, minimum width=52mm] (ask)
    {\textbf{Ask in the screen}\\{\scriptsize live or take-home? \dash{} is AI permitted?}};

  \node[kitbox, minimum width=34mm, below=10mm of ask, xshift=-40mm] (live)
    {\textbf{Live, AI permitted}\\{\scriptsize how you use it is scored}};
  \node[kitbox, minimum width=34mm, below=10mm of ask] (home)
    {\textbf{Take-home}\\{\scriptsize nobody is watching}};
  \node[kitbox, minimum width=34mm, below=10mm of ask, xshift=40mm] (free)
    {\textbf{AI not permitted}\\{\scriptsize fluency without it}};

  \draw[kitarrow] (ask) -- (live);
  \draw[kitarrow] (ask) -- (home);
  \draw[kitarrow] (ask) -- (free);

  \node[kitsoft, minimum width=112mm, below=9mm of home] (shared)
    {discovery \dash{} design \dash{} time-boxing \dash{} vertical slices \dash{} failure handling};

  \draw[kitarrow, draw=kitgrey] (live) -- (live |- shared.north);
  \draw[kitarrow, draw=kitgrey] (home) -- (shared);
  \draw[kitarrow, draw=kitgrey] (free) -- (free |- shared.north);

  \node[kitlabel, below=2mm of shared] {shared by all three \dash{} which is why preparing for all three is cheap};
\end{tikzpicture}
}{One question in the screening call decides which column you are in. Everything below the grey band is shared, which is why preparing for all three costs far less than three times as much.}

\section{What each format is really testing}

\textbf{The live build} tests how you behave under observation with an
unfinished problem. Not whether you can build the thing \dash{} they know you
can \dash{} but what you do when the clock is visible, what you cut, and
whether your reasoning is available to somebody watching.

\textbf{The take-home} tests what you do when nobody is watching: whether you
scope yourself, whether you test without being asked, and what you choose to
document. \reported{The defence conversation afterwards is reported to matter
more than the code alone.}{field-guide-take-home}

\textbf{The AI-free round} tests fluency, and it is the one where the recent
habit of reaching for a model is a genuine liability if you have not practised
without it.

\judgement{If you have limited rehearsal time, rehearse the live build. It is
the least forgiving of the three, its failure modes are the ones you cannot
recover from afterwards, and the discipline it demands \dash{} scoping out
loud, cutting deliberately, narrating decisions \dash{} improves the other two
as a side effect.}

\begin{drill}
For the process you are in, write down which format you are facing and what you
know about its rules. Every blank is a question for the coordinator. Then look
at the table above and mark the rows you have never practised \dash{} that is
your rehearsal list for Chapter~\ref{ch:rehearsal}.
\end{drill}

\section*{What this chapter does not claim}

The survey figure is one vendor's survey of hiring leaders and the sample size
was not available; the sentence citing it says it is one survey. The two
single-source claims are marked in the text and should be verified with the
company rather than relied on.

No claim is made about which way the split is moving. Both directions are
reported in the same corpus, and a confident prediction here would be exactly
the kind of claim this book refuses elsewhere.
